% Canonical equation -- KH instability numerical methods (KH-NM)
% kh-sim canonical-equations / v0.1
% Companion to kh-numerical-methods.md
% Bibliography: ../reference-bibliography/refs.bib
% Licence: CC-BY-SA-4.0

\documentclass[11pt,a4paper]{article}
\usepackage[utf8]{inputenc}
\usepackage[T1]{fontenc}
\usepackage{amsmath, amssymb, amsthm}
\usepackage{geometry}
\usepackage{booktabs}
\usepackage{hyperref}
\usepackage[backend=biber,style=numeric]{biblatex}
\addbibresource{../reference-bibliography/refs.bib}

\geometry{margin=2.4cm}

\title{Canonical equations:\\ Kelvin--Helmholtz numerical methods}
\author{kh-sim project}
\date{v0.1 -- 2026-05-12}

\theoremstyle{definition}
\newtheorem{definition}{Definition}
\theoremstyle{plain}
\newtheorem{proposition}{Proposition}

\begin{document}
\maketitle

This document specifies the numerical methods that the kh-sim solver
uses to discretise and integrate the governing equations stated in
the companion document \texttt{kh-governing-equations.tex}
(Eqs.~KH-GE-1..KH-GE-5). The numerical-methods specification is the
authoritative source for the \texttt{TC-NUM-KH-001..008} test cases
implemented in \texttt{backends/fortran/tests/test\_num\_*.f90}.

\section{Pseudo-spectral spatial discretisation (Eq.\ KH-NM-1)}

\begin{definition}[2D Fourier-spectral discretisation]
Discretise $\Omega = [0, L_x] \times [0, L_y]$ on a uniform grid of
$N_x \times N_y$ points. For a real-valued field $f$ on the grid,
the 2D real-input FFT (\texttt{rfft2}) produces a frequency-domain
representation $\hat f(k_x, k_y)$ on the half-grid
$N_x \times (N_y/2 + 1)$, where the wavenumber arrays are
\begin{align*}
k_x &= \frac{2\pi}{L_x}\,\{0, 1, 2, \ldots, \tfrac{N_x}{2}-1, -\tfrac{N_x}{2}, \ldots, -1\},\\
k_y &= \frac{2\pi}{L_y}\,\{0, 1, 2, \ldots, \tfrac{N_y}{2}\}.
\end{align*}
Spatial derivatives are diagonal in this basis:
\begin{equation}
\label{eq:khnm1}
\widehat{\partial_x f} = i k_x \hat f,
\qquad
\widehat{\partial_y f} = i k_y \hat f,
\qquad
\widehat{\nabla^2 f} = -(k_x^2 + k_y^2)\,\hat f.
\tag{KH-NM-1}
\end{equation}
\end{definition}

\section{Spectral Poisson solver (Eq.\ KH-NM-2)}

\begin{proposition}[Diagonal Poisson]
The streamfunction Poisson equation
$\partial_{xx}\psi + \partial_{yy}\psi = -\omega$
(Eq.~KH-GE-2) is diagonal in the Fourier basis of \eqref{eq:khnm1}.
The solution is
\begin{equation}
\label{eq:khnm2}
\hat\psi(k_x, k_y) = \frac{\hat\omega(k_x, k_y)}{k_x^2 + k_y^2}
\quad\text{for}\quad (k_x, k_y) \neq (0, 0),
\qquad
\hat\psi(0, 0) = 0.
\tag{KH-NM-2}
\end{equation}
\end{proposition}

\begin{proof}
Apply $\widehat{\nabla^2 f} = -(k_x^2 + k_y^2)\hat f$ to both sides
of KH-GE-2:
\[
-(k_x^2 + k_y^2)\hat\psi = -\hat\omega,
\]
which rearranges directly. The choice $\hat\psi(0,0) = 0$ fixes the
gauge freedom in $\psi$ (the Poisson equation determines $\psi$ only
up to an additive constant).
\qed
\end{proof}

This corresponds to test case \texttt{TC-NUM-KH-002}.

\section{Time integration}

\subsection{Canonical Runge--Kutta 4 (Eq.\ KH-NM-3a)}

Writing the vorticity-transport equation KH-GE-1 abstractly as
$\partial_t \omega = f(\omega)$ with
\[
f(\omega) = -u\,\partial_x \omega - v\,\partial_y \omega + \nu\,\nabla^2 \omega,
\]
where $u, v$ are recovered from $\omega$ via \eqref{eq:khnm2} +
KH-GE-3, the classical 4th-order Runge--Kutta scheme advances
$\omega$ from $t^n$ to $t^{n+1} = t^n + \Delta t$ via
\begin{equation}
\label{eq:khnm3a}
\begin{aligned}
k_1 &= f(\omega^n),\\
k_2 &= f(\omega^n + \tfrac{1}{2}\Delta t\,k_1),\\
k_3 &= f(\omega^n + \tfrac{1}{2}\Delta t\,k_2),\\
k_4 &= f(\omega^n + \Delta t\,k_3),\\
\omega^{n+1} &= \omega^n + \tfrac{\Delta t}{6}\,(k_1 + 2k_2 + 2k_3 + k_4).
\end{aligned}
\tag{KH-NM-3a}
\end{equation}
RK4 has $\mathcal{O}(\Delta t^5)$ local truncation error and
$\mathcal{O}(\Delta t^4)$ global error in the well-resolved regime.
This is the canonical scheme used by \texttt{shared/physics/kh\_physics.py}
and \texttt{backends/rust/src/physics/solver.rs}. The standard reference
is Hairer, N{\o}rsett \& Wanner~\cite{HairerNorsettWanner1993}.

\subsection{Alternate ETDRK4 (Eq.\ KH-NM-3b)}

When the viscous diffusion term $\nu\nabla^2\omega$ is stiff (high Re
or under-resolved grids), the \emph{exponential time-differencing
Runge--Kutta 4} scheme of Cox \& Matthews~\cite{CoxMatthews2002}
handles the linear part exactly via matrix exponentiation, with
RK4-like stages for the non-linear part. Split the equation as
\[
\partial_t \omega = L\,\omega + N(\omega),
\]
where $L = \nu\nabla^2$ is diagonal in Fourier space
($L \leftrightarrow -\nu(k_x^2 + k_y^2)$) and
$N(\omega) = -u\,\partial_x\omega - v\,\partial_y\omega$ is the
non-linear advection. The ETDRK4 update is
\begin{equation}
\label{eq:khnm3b}
\omega^{n+1} = e^{L\Delta t}\,\omega^n
+ \Delta t\,\bigl[\phi_1 N_1 + 2\phi_2(N_2 + N_3) + \phi_3 N_4\bigr]
\tag{KH-NM-3b}
\end{equation}
where the $\phi_i$ are the standard $\phi$-functions
(see Kassam \& Trefethen~\cite{KassamTrefethen2005} for the
contour-integral-evaluated coefficients that stabilise the small-$L$
limit), and $N_i$ are non-linear-term evaluations at intermediate
stages with RK4-like sub-step structure. ETDRK4 is the scheme used by
\texttt{backends/fortran/src/kh\_etdrk4.f90}.

The two schemes \eqref{eq:khnm3a} and \eqref{eq:khnm3b} produce
identical solutions in the limit of well-resolved grids and modest
Reynolds numbers, but diverge in stiff regimes. The kh-sim project
flags this as a known disambiguation question
(AUDIT-kh-sim-2026-05-12 §6.1). Test cases
\texttt{TC-NUM-KH-003} (ETDRK4 on linear scalar ODE) and
\texttt{TC-NUM-KH-004} (temporal convergence order $4$) anchor the
ETDRK4 implementation.

\section{Aliasing control: 2/3-rule (Eq.\ KH-NM-4)}

\begin{definition}[2/3-rule dealiasing]
After every evaluation of a non-linear product (e.g.
$u \cdot \partial_x \omega$) in physical space and subsequent
forward FFT, set
\begin{equation}
\label{eq:khnm4}
\hat g(k_x, k_y) := 0
\quad\text{whenever}\quad
\lvert k_x \rvert > \frac{N_x}{3} \;\text{or}\; \lvert k_y \rvert > \frac{N_y}{3}.
\tag{KH-NM-4}
\end{equation}
\end{definition}

The surviving Fourier modes contain only the aliasing-free
representation of the product. This is the standard 2/3-rule of
Orszag~\cite{Orszag1971}.
This corresponds to test case \texttt{TC-NUM-KH-005}.

\section{Stability constraint: CFL (Eq.\ KH-NM-5)}

\begin{proposition}[CFL stability for explicit RK4 + Fourier spectral]
For the explicit RK4 time integrator
\eqref{eq:khnm3a} on a periodic Fourier-spectral spatial
discretisation \eqref{eq:khnm1}, the maximum stable timestep
satisfies
\begin{equation}
\label{eq:khnm5}
\Delta t \leq \min\!\left(\frac{\Delta x}{\lvert u \rvert_{\max}},\,
                          \frac{\Delta y}{\lvert v \rvert_{\max}},\,
                          \frac{(\Delta x)^2}{2\nu},\,
                          \frac{(\Delta y)^2}{2\nu}\right).
\tag{KH-NM-5}
\end{equation}
\end{proposition}

The kh-sim default $\Delta t = 0.001$ is conservative for the
default $\mathrm{Re} = 1000$ on a $128 \times 64$ grid. This
corresponds to test case \texttt{TC-NUM-KH-006}. The ETDRK4 scheme
\eqref{eq:khnm3b} relaxes the diffusive bound (last two min-arguments)
by treating $L$ exactly; the advective bound remains.

\section{Numerical-precision considerations}

By IEEE 754-2019~\cite{IEEE754_2019}, the rounding bound for the
inner-product operations in the RK4 stages at IEEE binary64 (double
precision) is approximately $\varepsilon \approx 2.22 \times 10^{-16}$
per operation. For an integration of $N_t$ steps on an $N_x \times N_y$
grid with the non-linear advection involving $\mathcal{O}(N_x N_y)$
operations per RHS evaluation, the cumulative round-off bound is
\[
\varepsilon_{\text{cum}}(N_t, N_x, N_y)
\;\lesssim\;
\mathcal{O}\bigl( N_t \cdot \sqrt{N_x N_y} \cdot \varepsilon \bigr)
\quad\text{(benign-input statistical bound)},
\]
with the worst-case linear-in-$N$ bound rarely realised. For the
kh-sim default reference run ($N_t = 100$, $N_x N_y = 8192$ at
$64 \times 128$ default), this gives $\varepsilon_{\text{cum}} \sim
10^{-12}$, an order of magnitude tighter than the $\pm 5\%$ diagnostic
gate currently used (AUDIT-kh-sim §4 lists this as an open question
for phase 5 gate derivation).

\section{Test-case traceability}

\begin{table}[h]
\centering
\caption{Mapping from KH-NM equations to TC-NUM-KH test cases.}
\begin{tabular}{lll}
\toprule
\textbf{Spec} & \textbf{KH-NM section} & \textbf{Fortran test} \\
\midrule
TC-NUM-KH-001 & KH-NM-1 (2D FFT round-trip)            & \texttt{test\_num\_001\_fft\_roundtrip.f90} \\
TC-NUM-KH-002 & KH-NM-2 (spectral Poisson)              & \texttt{test\_num\_002\_poisson.f90} \\
TC-NUM-KH-003 & KH-NM-3b (ETDRK4 on linear scalar ODE) & \texttt{test\_num\_003\_etdrk4\_linear.f90} \\
TC-NUM-KH-004 & KH-NM-3b (convergence order 4)          & \texttt{test\_num\_004\_convergence.f90} \\
TC-NUM-KH-005 & KH-NM-4 (2/3-rule)                       & \texttt{test\_num\_005\_dealias.f90} \\
TC-NUM-KH-006 & KH-NM-5 (CFL across Re)                  & \texttt{test\_num\_006\_cfl.f90} \\
TC-NUM-KH-007 & KH-NM-3a + KH-NM-2 (energy conservation, inviscid) & \texttt{test\_num\_007\_energy.f90} \\
TC-NUM-KH-008 & KH-NM-1..5 composed (reference run validation) & \texttt{test\_num\_008\_reference.f90} \\
\bottomrule
\end{tabular}
\end{table}

\printbibliography

\end{document}
